% ===================================================================
% 03_device_data.tex -- Device, measured data and inputs (writer C1). ASCII only.
% ===================================================================
\chapter{Device, measured data and inputs}\label{ch:device}

A TCAD model cannot be more specific than the information on which it is built. This chapter therefore states, item by
item, what is known about the measured transistors, where each item comes from, and what is not known. It addresses four
questions: (i) which device and process are modelled (\cref{sec:inp-device}); (ii) how the physical stack is translated
into the two-dimensional simulation domain (\cref{sec:inp-geometry}); (iii) what the measured data consist of
(\cref{sec:inp-data}) and which figures of merit they yield when the extractor used for the simulations is applied to
them (\cref{sec:inp-metrics-measured}); and (iv) which gaps in the data limit every later conclusion
(\cref{sec:inp-limitations}). No quantity in this chapter is fitted. Every number is either measured on the target
devices, read from the user's schematic, or explicitly labelled as assumed.

% -------------------------------------------------------------------
\section{Device and process}\label{sec:inp-device}

The devices are the bottom-gate IWO thin-film transistors of the thickness series of \citet{Januar2026}. The primary
source for the stack is Sec.~2.1 of that paper \citep[pp.~2--3]{Januar2026} together with the user's device schematic
\citep{DeviceSchematic}; the measured transfer curves are taken from the experimental workbook
\citep{ExperimentalWorkbook}.

From the substrate upwards, the stack consists of an n$^+$ Si support carrying a \SI{50}{nm} TiN bottom gate deposited
by physical vapour deposition. On the gate lie a \HfO{} gate dielectric (nominal \SI{15}{nm}) and an \AlO{} interlayer
(nominal \SI{2}{nm}), both grown by ALD at \SI{250}{\celsius}. The paper reports growth rates of about 0.92 and
\SI{1.04}{\angstrom} per cycle for the two oxides \citep[p.~2]{Januar2026}, while the nominal thicknesses of 15 and
\SI{2}{nm} are taken from the schematic; the \AlO{} layer is in contact with the channel. The channel is a sputtered IWO
film (about 2\,\% W; the definition of the composition cannot be determined from the available data) of nominal
thickness 2.0, 6.3, 13.2 or \SI{31.8}{nm}. \SI{70}{nm} Pd source and drain electrodes are placed on top of the channel
and patterned with contact holes and metallization. There is no passivation or capping layer, so the channel between
source and drain is exposed to the ambient.

A rapid thermal anneal at \SI{150}{\celsius} for \SI{600}{s} completes the process (process record of the package,
\protect\file{docs/DEVICE_STRUCTURE.md}); it serves as process context only and does not enter any model parameter
directly. The channel length and width, $L = \SI{20}{\micro\metre}$ and $W = \SI{290}{\micro\metre}$, and the drain bias
of the transfer measurement, $\Vd = \SI{0.7}{V}$, are read from the schematic, because the workbook columns carry no
units.

\Cref{tab:inp-stack} lists every structural item with its provenance label. The \HfO{} permittivity 19.57 and the IWO
permittivity 9.30 are C-V results reported in the SI of \citet{Januar2026} (Fig.~S1; the SI is not bundled, and the
values were recorded by the Astra audit; the main text gives only ``approximately 9'' for IWO,
\citealp[p.~2]{Januar2026}). The \AlO{} permittivity 9.0 is assumed. The consequences of each label for the gate
capacitance are derived in \cref{sec:par-stack}.

\begin{table}[htbp]
\centering
\small
\caption{Structural items of the modelled device and their provenance (\protect\file{config/geometry.yaml},
\protect\file{docs/DEVICE_STRUCTURE.md}).}
\label{tab:inp-stack}
\begin{tabularx}{\linewidth}{@{}lXl@{}}
\toprule
Item & Value & Provenance \\
\midrule
Gate & TiN \SI{50}{nm} on n$^+$ Si (Si screened by the continuous TiN, omitted) & \prov{MEASURED\_TARGET\_DEVICE} \\
Gate dielectric & \HfO{} \SI{15.0}{nm}, $\varepsilon_r = 19.57$ (SI Fig.~S1) & \prov{MEASURED\_TARGET\_DEVICE} \\
Interlayer & \AlO{} \SI{2.0}{nm}, $\varepsilon_r = 9.0$ & thickness \prov{MEASURED\_TARGET\_DEVICE}; $\varepsilon_r$ \prov{ASSUMED} \\
Channel & IWO, $t = 2.0/6.3/13.2/31.8$\,nm, $\epsIWO = 9.3$ & \prov{MEASURED\_TARGET\_DEVICE} \\
Source/drain & Pd \SI{70}{nm}, top contacts & \prov{MEASURED\_TARGET\_DEVICE} \\
S/D overlap & \SI{2}{\micro\metre} per side & \prov{ASSUMED} (not reported) \\
Top of channel & air, no cap & as reported (no fictitious cap added) \\
$W/L$ & \SI{290}{\micro\metre} / \SI{20}{\micro\metre} & \prov{MEASURED\_TARGET\_DEVICE} (schematic) \\
$\Vd$ (transfer) & \SI{0.7}{V} & \prov{MEASURED\_TARGET\_DEVICE} (schematic) \\
Temperature & \SI{300}{K} & \prov{ASSUMED} (not recorded) \\
Current unit & A/$\mu$m & interpretation of the schematic axis \\
\bottomrule
\end{tabularx}
\end{table}

\begin{figure}[htbp]
\centering
\includegraphics[width=\linewidth]{stack_schematic}
\caption{Cross-section of the modelled device (not to scale): TiN gate \SI{50}{nm}, \HfO{} \SI{15}{nm}, \AlO{}
\SI{2}{nm}, IWO channel of thickness $t$, Pd source/drain \SI{70}{nm} overlapping the channel by \SI{2}{\micro\metre}
on each side, air above the channel; $L = \SI{20}{\micro\metre}$, $W = \SI{290}{\micro\metre}$. The ATLAS region
numbers of the physical-structure deck are indicated (1: IWO bulk; 2: \SI{0.25}{nm} IWO interface layer; 3: \AlO{};
4: \HfO{}; 6: Conductor gate volume; 7: Air). Drawn from \protect\file{config/geometry.yaml} and
\protect\file{docs/DEVICE_STRUCTURE.md}.}
\label{fig:stack_schematic}
\end{figure}

% -------------------------------------------------------------------
\section{Geometry and the simulated structure}\label{sec:inp-geometry}

\paragraph{Coordinates.} The simulation domain is a two-dimensional cross-section along the channel. In micrometres,
$x$ runs from 0 to 24: the source contact occupies $0 \le x \le 2$, the channel $2 < x < 22$ and the drain
$22 \le x \le 24$. The vertical coordinate $y$ increases downward from the IWO/\AlO{} interface at $y = 0$. The IWO film
occupies $-t \le y \le 0$, the \AlO{} $0 \le y \le 0.002$, the \HfO{} $0.002 \le y \le 0.017$ and the TiN gate volume
$0.017 \le y \le 0.067$; the Pd electrodes and the air region occupy $-t-0.07 \le y \le -t$. The \SI{2}{\micro\metre}
overlap of the contacts is not reported. With contacts on the top surface of a \SI{20}{\micro\metre} long channel, the
drain current is expected to be insensitive to it. However, the one launch intended to test this (\run{0027}, overlap
\SI{4}{\micro\metre}) is malformed, and its result is retracted (\cref{sec:cal-param-control}).

\paragraph{Width convention.} The mesh is given the width \SI{1}{\micro\metre} (\texttt{mesh width=1.0}), so that every
logged current is numerically in A/$\mu$m and can be compared directly with the workbook. The current of the physical
device is $290 \times$ the logged value; a division by 290 is never applied. This convention was verified in the V0
package (V0 run\_0007 with \texttt{WIDTH=2} against V0 run\_0005: raw currents exactly $\times 2.000000$, currents per
micrometre identical; check J of \cref{tab:num-convergence}).

\paragraph{Two electrically equivalent representations.} The deck generator \protect\file{scripts/build_iwo_decks.py}
writes the stack in one of two forms (\cref{sec:num-deck}):
\begin{enumerate}
  \item \emph{Reduced structure} (runs \run{0001}--\run{0009}): IWO regions, \AlO{} ($\varepsilon_r = 9.0$) and \HfO{}
        ($\varepsilon_r = 19.57$). The gate is a Dirichlet boundary at the bottom of the \HfO{} with the TiN work
        function, and source and drain are Dirichlet boundaries on the IWO top surface over $0$--$2$ and
        $22$--$24$\,$\mu$m. In DC a continuous metal is an equipotential, so replacing its volume by a boundary with the
        same work function changes neither the field in the dielectrics nor the contact boundary condition.
  \item \emph{Physical structure} (\texttt{structure: full}, every run from \run{0012} on): the same IWO regions;
        explicit \AlO{} and \HfO{} regions whose permittivities are set on \texttt{material region=} statements
        (the ATLAS defaults for these material names are not the measured values); a \SI{50}{nm}
        \texttt{Conductor} volume carrying the gate electrode (TiN is not in the ATLAS metal table; work function
        \SI{4.70}{eV} set on \texttt{contact}, \prov{ASSUMED}); a \SI{70}{nm} Air region above the channel; and
        \SI{70}{nm} $\times$ \SI{2}{\micro\metre} Palladium source/drain electrode volumes without a work function
        (Ohmic, \prov{ASSUMED}; no TLM data).
\end{enumerate}
The equivalence of the two forms was verified in the real solver. \run{0012} (physical structure, stepped sweep) and
\run{0008} (reduced structure, pointwise sweep), with identical parameters, differ by $+0.08$, $+0.05$, $+0.01$,
$-0.003$ and $-0.002$\,\% in $\Id$ at $\Vg = 0.5$, 0.7, 1, 2 and \SI{3}{V} (check M in \cref{tab:num-convergence}).

\paragraph{Gate capacitance.} With $\varepsilon_0 = \SI{8.854e-14}{\farad\per\centi\metre}$, the series capacitance of
the two oxides is
\begin{equation}
  \Cox = \frac{\varepsilon_0}{t_{\mathrm{HfO_2}}/\varepsilon_{\mathrm{HfO_2}} + t_{\mathrm{Al_2O_3}}/\varepsilon_{\mathrm{Al_2O_3}}}
       = \frac{8.854\times10^{-14}}{(1.5\times10^{-6}/19.57) + (2\times10^{-7}/9.0)}\ \mathrm{F\,cm^{-2}}
       = 8.955\times10^{-7}\ \mathrm{F\,cm^{-2}},
  \label{eq:inp-cox}
\end{equation}
which is equivalent to $\EOT = 3.9\,\varepsilon_0/\Cox = \SI{3.86}{nm}$. A sheet charge of $10^{12}\cmsq$ corresponds to
$q\times10^{12}/\Cox = \SI{0.179}{V}$. These three numbers are used throughout the report; their derivation and error
budget are given in \cref{sec:par-stack}.

% -------------------------------------------------------------------
\section{The measured data}\label{sec:inp-data}

The measurement consists of a single transfer curve per thickness, stored in the package as
\protect\file{data/experimental_clean.csv} with the columns \texttt{thickness\_nm}, \texttt{source\_row}, \texttt{vg\_V}
and \texttt{id\_A\_per\_um}. Each curve spans $-3 \le \Vg \le 3$\,V in steps of \SI{0.05}{V} at $\Vd = \SI{0.7}{V}$.
No other electrical data exist for these devices (\cref{sec:inp-limitations}).

\begin{figure}[htbp]
\centering
\includegraphics[width=\linewidth]{measured_transfer}
\caption{Measured transfer curves of the 2.0, 6.3, 13.2 and \SI{31.8}{nm} devices at $\Vd = \SI{0.7}{V}$
(\protect\file{data/experimental_clean.csv}). (a) Logarithmic scale; the dotted lines mark the off-band medians
($-2 \le \Vg \le -0.5$\,V): $4.61\times10^{-15}$ (2.0 nm), $1.53\times10^{-13}$ (6.3 nm) and
$6.60\times10^{-12}\Aum$ (13.2 nm). (b) Linear scale. The 2.0 and \SI{6.3}{nm} curves are close to each other, the
\SI{13.2}{nm} curve is shifted to lower gate voltage and carries about six times more on-current, and the
\SI{31.8}{nm} device does not turn off within the sweep.}
\label{fig:measured_transfer}
\end{figure}

Four features of \cref{fig:measured_transfer} determine the course of the whole study.
\begin{enumerate}
  \item \emph{Non-monotonic thickness dependence.} The \SI{2.0}{nm} and \SI{6.3}{nm} curves nearly coincide in
        threshold and on-current, whereas the \SI{13.2}{nm} curve differs strongly; the \SI{6.3}{nm} film even has
        the lower field-effect mobility of the two thin films (\cref{sec:inp-metrics-measured}). A smooth thickness
        law cannot describe this behaviour without a device-specific parameter (\cref{ch:6p3}).
  \item \emph{Off-state floors.} Each curve saturates at a flat plateau in the off-band. The plateaus increase steeply
        with thickness: the \SI{13.2}{nm} floor is 1432 times the \SI{2}{nm} floor, which corresponds to a log-log
        slope of 3.78 over the three thin films (the S4 report quotes 3.85; \cref{fig:off_floors}). Drift-diffusion
        with thermal generation does not produce such floors; the Astra audit bounded the thermal-generation current
        at 13 to 17 decades below them. Gate leakage is disfavoured as the dominant floor at 6.3 and \SI{13.2}{nm},
        because the floors do not follow the drain--gate bias (\cref{sec:mech-traps}). Their origin (a measurement
        floor or an ungated channel path) cannot be determined from the available data, because $\Ig$ and $\Is$ were
        not recorded.
  \item \emph{An isolated low point.} The minimum of the \SI{2}{nm} curve, about $1.1\times10^{-16}\Aum$ at
        $-3$\,V, is a single point about 40 times below the floor of that curve. It is treated as an outlier, and the
        KCL acceptance limit of the runner was redefined on the median off-band current instead of the minimum
        (\cref{sec:num-validation}).
  \item \emph{The \SI{31.8}{nm} device is always on} ($\Id \ge 2.6\times10^{-7}\Aum$ at $\Vg = -3$\,V). It lies
        outside the \mbox{2--13\,nm} series published by \citet[p.~2]{Januar2026}, was excluded by the user from the
        final calibration set, and appears only in the early runs \run{0006} and \run{0010}. The V0 package required
        an additional back-surface donor sheet and a series resistance of about $9\times10^{4}\,\Omega\,\mu$m to
        describe it.
\end{enumerate}

% -------------------------------------------------------------------
\section{Measured metrics}\label{sec:inp-metrics-measured}

\Cref{tab:measured-metrics} lists the figures of merit of the measured curves. They are extracted by
\protect\file{scripts/extract_metrics.py}, the same code that scores every ATLAS run, so that measured and simulated
numbers are directly comparable; the definitions and their derivations are given in \cref{sec:der-metrics}. The swing is
reported over fixed current windows, because the minimum local swing $\SSmin$ depends on the voltage grid and on the
floor window and is therefore not a reliable metric (\cref{sec:der-metrics}); it is listed only for comparison with
earlier documents.

\begin{table}[htbp]
\centering
\small
\caption{Measured metrics at $\Vd = \SI{0.7}{V}$ (\evDATA; \protect\file{data/experimental_clean.csv}, extractor
\protect\file{scripts/extract_metrics.py}; numbers from \protect\file{figures/MANIFEST.md} and the evidence brief of 2026-09-25).
Currents in A/$\mu$m, swings in mV/dec, mobilities in \si{cm^2/(V.s)}.}
\label{tab:measured-metrics}
\begin{tabular}{@{}lrrrr@{}}
\toprule
Metric & 2.0 nm & 6.3 nm & 13.2 nm & 31.8 nm \\
\midrule
$\Vthcc$ at $10^{-9}\Aum$ (V) & 0.662 & 0.644 & 0.083 & always on \\
$\Vthlin$ (V) & 1.663 & 1.820 & 1.044 & -- \\
$\Vthlin - \Vthcc$ (V) & 1.001 & 1.176 & 0.962 & -- \\
$\SSfix{10^{-11}}{10^{-10}}$ & 121.2 & 124.5 & 136.1 & -- \\
$\SSfix{10^{-10}}{10^{-9}}$ & 194.8 & 210.0 & 121.6 & -- \\
$\SScc = \SSfix{10^{-10}}{10^{-8}}$ & 269.9 & 295.4 & 160.0 & -- \\
$\SSmin$ (not robust) & 84.5 & 114.7 & 130.8 & -- \\
$\gmmax$ (A/(V\,$\mu$m)) & $3.807\times10^{-7}$ & $3.431\times10^{-7}$ & $1.572\times10^{-6}$ & -- \\
$\muFE$ (linear, $\gm$-based) & 12.15 & 10.95 & 50.17 & -- \\
$\musat$ (saturation formula) & 5.09 & 3.86 & 27.40 & -- \\
$\Ion = \Id(\Vg = \SI{3}{V})$ & $5.09\times10^{-7}$ & $4.03\times10^{-7}$ & $3.07\times10^{-6}$ & $3.59\times10^{-6}$ \\
Off-floor (median, $-2$ to $-0.5$\,V) & $4.61\times10^{-15}$ & $1.53\times10^{-13}$ & $6.60\times10^{-12}$ & -- \\
\bottomrule
\end{tabular}
\end{table}

\begin{figure}[htbp]
\centering
\includegraphics[width=\linewidth]{measured_metrics}
\caption{Measured metrics versus channel thickness (2.0, 6.3, \SI{13.2}{nm}; \SI{31.8}{nm} as a hollow symbol where
defined): $\Vthcc$ and $\Vthlin$, fixed-current swings $\SSfix{10^{-11}}{10^{-10}}$ and $\SSfix{10^{-10}}{10^{-8}}$,
linear $\muFE$ and saturation-formula mobility, $\Ion$ and the off-floor. Same numbers as
\cref{tab:measured-metrics}. The step between 6.3 and \SI{13.2}{nm} in threshold, mobility and on-current, and the
relation $\muFE(6.3) < \muFE(2.0)$, are visible.}
\label{fig:measured_metrics}
\end{figure}

\paragraph{Interpretation of the table.} Three observations are essential for the later chapters. (i) $\Vthcc$, $\muFE$
and $\Ion$ are non-monotonic or step-like: the two thinnest films are similar, whereas the \SI{13.2}{nm} film differs by
about \SI{0.58}{V} in threshold and by a factor of about four in mobility. (ii) The \SI{6.3}{nm} film has the largest gap
$\Vthlin - \Vthcc$ (\SI{1.176}{V}); this on-state shape is the quantity that no parameter set of the model reproduces
(\cref{ch:6p3}). (iii) The swing in the window $10^{-10}$--$10^{-8}\Aum$ is 270--\SI{295}{mV/dec} for the two thin films
and \SI{160}{mV/dec} for the \SI{13.2}{nm} film. These values are far above the thermal limit of about \SI{60}{mV/dec}
at \SI{300}{K}, which indicates a large density of states in the gap.

\paragraph{Mobility definitions and the values of the paper.} \citet{Januar2026} quote peak mobilities of about 5.1 and
\SI{27.4}{cm^2/(V.s)} for the thinnest and thickest films of the series \citep[p.~3]{Januar2026} and relate the
decrease from about 27.4 to about 5.1 to a roughly fourfold increase of the trap density \citep[p.~7]{Januar2026}. The
linear, transconductance-based $\muFE$ of the workbook curves is 12.15 and \SI{50.17}{cm^2/(V.s)}, i.e.\ a factor of
2.4 and 1.8 higher. Specialist S5 showed that applying the saturation formula to the same curves at
$\Vd = \SI{0.7}{V}$ gives 5.09 and \SI{27.40}{cm^2/(V.s)}, which reproduces the numbers of the paper. The workbook
curves are therefore the thickness series of the paper, and the mobility of the paper is a saturation-formula peak
taken near threshold rather than a linear-regime field-effect mobility (\evCORR; \cref{sec:der-mufe},
\cref{fig:mufe_extraction}). The model is calibrated against the curves themselves and never against a mobility number,
so the choice of definition affects the comparison with the literature but not the calibration.

% -------------------------------------------------------------------
\section{Data limitations and their consequences}\label{sec:inp-limitations}

The following items are not available. Each limitation is stated together with the conclusion that it prevents.

\begin{enumerate}
  \item One device per thickness. The number of devices per thickness and their spread are unknown. As a consequence,
        a difference between two films that is smaller than the (unknown) device-to-device spread cannot be
        attributed to thickness; in particular, the \SI{6.3}{nm} anomaly may be a property of the device rather than
        of the thickness.
  \item No gate or source current. The off-state floors therefore cannot be assigned directly to an instrument floor
        or to channel conduction (gate leakage is only disfavoured indirectly, from the bias independence of the
        floors), and the KCL check cannot be compared with the measurement.
  \item No sweep direction, hysteresis or dwell time. Slow traps and bias-stress shifts therefore cannot be separated
        from the static DOS, and the fitted $\Qf$ may absorb an offset caused by the sweep history.
  \item No output ($\Id$--$\Vd$) curves. Series resistance and contact behaviour are therefore under-constrained, and
        the $\Id$--$\Vd$ curves of \cref{ch:predictions} are predictions, not fits.
  \item No C-V on these devices and no TLM. $\Cox$ therefore relies on the SI permittivity of \HfO{} and an assumed
        \AlO{} permittivity, and the Ohmic contact assumption cannot be tested directly (\cref{sec:par-contacts}).
  \item No temperature record. The measurement temperature is assumed to be \SI{300}{K}, so the thermal voltage in all
        swing and occupancy arguments carries this assumption. Temperature-dependent data exist in the SI of
        \citet[SI Fig.~S12, caption on p.~11]{Januar2026} (2 nm at 338 and \SI{358}{K}) but are not available to the
        project, which is why the temperature runs of \cref{ch:predictions} are registered as predictions.
  \item No thickness uncertainty. The method used to measure the film thickness and its uncertainty are not reported.
        A thickness error of the order of a nanometre at \SI{6.3}{nm} is therefore an open hypothesis (tested as A1,
        \run{0031}, in \cref{ch:6p3}).
  \item No ambient record. The channel is uncapped, and adsorbates on the exposed back surface can carry charge. A
        back-surface charge $\Qb$ therefore cannot be excluded and is tested as a hypothesis (A3, \run{0030}).
  \item Units inferred from the schematic. The workbook headers carry no units; A/$\mu$m and $\Vd = \SI{0.7}{V}$ are
        taken from the schematic. The consistency of the saturation-formula mobility with the paper
        (\cref{sec:inp-metrics-measured}) supports this reading.
  \item Composition. ``About 2\,\% W'' is not defined as an atomic or a weight fraction in the available sources. No
        parameter depends on it directly, and $\Ndeff$ is an electrically active donor density, not the W
        concentration (\cref{sec:par-nd}).
  \item The \SI{31.8}{nm} column. It lies outside the published series and cannot be turned off within the sweep; it is
        excluded from the final calibration set by the user's decision.
\end{enumerate}

\begin{keyresult}
The model rests on one measured transfer curve per thickness (\evDATA), on a gate stack whose capacitance
$\Cox = \SI{8.955e-7}{\farad\per\centi\metre\squared}$ ($\EOT = \SI{3.86}{nm}$) follows from a measured \HfO{} and an
assumed \AlO{} permittivity, and on an assumed Ohmic top-contact geometry. The measured metrics are non-monotonic in
thickness: the 2.0 and \SI{6.3}{nm} films are similar ($\Vthcc$ 0.662 and \SI{0.644}{V}, $\muFE$ 12.15 and
\SI{10.95}{cm^2/(V.s)}), whereas the \SI{13.2}{nm} film is very different ($\Vthcc = \SI{0.083}{V}$,
$\muFE = \SI{50.17}{cm^2/(V.s)}$). The mobilities of the paper are saturation-formula values (5.09 and 27.40 are
reproduced). The absence of $\Ig$/$\Is$, hysteresis, $\Id$--$\Vd$, C-V, TLM, temperature and device statistics limits
every later conclusion.
\end{keyresult}

The next two chapters derive, one input at a time, how this device and these data are translated into the parameters of
the ATLAS model, beginning with the electrostatic inputs.
